%%%%%%%%%%%%%%%%%%%%%%%%%%%%%%%%%%%%%%%%%
% Lachaise Assignment
% LaTeX Template
% Version 1.0 (26/6/2018)
%
% This template originates from:
% http://www.LaTeXTemplates.com
%
% Authors:
% Marion Lachaise & François Févotte
% Vel (vel@LaTeXTemplates.com)
%
% License:
% CC BY-NC-SA 3.0 (http://creativecommons.org/licenses/by-nc-sa/3.0/)
% 
%%%%%%%%%%%%%%%%%%%%%%%%%%%%%%%%%%%%%%%%%

%----------------------------------------------------------------------------------------
%	PACKAGES AND OTHER DOCUMENT CONFIGURATIONS
%----------------------------------------------------------------------------------------

\documentclass{article}

\input{structure.tex} % Include the file specifying the document structure and custom commands

%----------------------------------------------------------------------------------------
%	ASSIGNMENT INFORMATION
%----------------------------------------------------------------------------------------

\title{Metodologías-2025: Evaluación Integral \#2} % Title of the assignment

\author{Luis Fernando González Avila} % Author name and email address

\date{Consejo Latinoamericano de Ciencias Sociales --- \today} % University, school and/or department name(s) and a date

%----------------------------------------------------------------------------------------

\begin{document}

\maketitle % Print the title

%----------------------------------------------------------------------------------------
%	INTRODUCTION
%----------------------------------------------------------------------------------------
\section*{Resumen Proporcionado}

\begin{scriptsize}
	
Las condiciones laborales precarias de las trabajadoras domésticas constituyen un fenómeno conocido. Sin embargo, desde hace más de una década se registran diferentes iniciativas, tanto a nivel nacional como internacional, tendientes a equiparar los derechos laborales de este colectivo con los de los demás asalariados. En este contexto se registra un renovado interés en materia de producción académica sobre el sector. No obstante, el tema de la segmentación de esta fuerza laboral, y sus implicancias en el impacto de estas iniciativas, ha sido todavía poco estudiado. En este sentido, una de las variables más significativas a la hora de considerar la segmentación de este colectivo se vincula con las modalidades de inserción diferenciales en la ocupación - fundamentalmente relacionadas con la dedicación horaria- que implican situaciones disímiles en términos de las experiencias ocupacionales y el acceso a derechos laborales. En este artículo nos proponemos contribuir a esta línea de indagación en base a una primera clasificación del universo de estas trabajadoras en dos subgrupos: aquéllas que se insertan pocas horas semanales bajo la modalidad de pago “por horas” y el resto de las trabajadoras, que se desempeña por jornadas extendidas bajo la modalidad de pago mensual. Combinando datos cuantitativos y cualitativos, se indagará, en primer lugar, sobre las motivaciones y/o restricciones que llevan a las trabajadoras a cada tipo de inserción. Se estudiará, luego, las diferentes formas en que cada subgrupo concibe y gestiona la relación con los empleadores, así como sus percepciones y apropiaciones de los derechos laborales.\\

\noindent\textit{Pereyra, Francisca; Tizziani, Ania; Experiencias y condiciones de trabajo diferenciadas en el servicio doméstico : hacia una caracterización de la segmentación laboral del sector en la ciudad de Buenos Aires ; Universidad Nacional de Santiago del Estero. Facultad de Humanidades y Ciencias Sociales. Programa de Investigaciones sobre Trabajo y Sociedad; Trabajo y sociedad; 23; 12-2014; 5-25}
	
\end{scriptsize}


\section{Definir claramente el objeto de estudio} % Unnumbered section

El objetivo del estudio es \textbf{analizar cómo las distintas modalidades de inserción laboral en el trabajo doméstico remunerado \textit{—particularmente la cantidad de horas trabajadas y el tipo de remuneración—} influyen en la relación que las trabajadoras establecen con sus empleadores y en su apropiación y acceso efectivo a los derechos laborales, considerando tanto factores estructurales como percepciones subjetivas}.\\

\noindent Este objetivo se alinea con lo que plantea Piovani et al.,2007 sobre una necesidad en la que los objetivos son la guía estratégica del diseño, orientando dede el inicio la división del objeto de estudio y la selección de técnicas. La anterior definición también reitera el principio de relevancia social del conocimiento al abordar una problemática concreta con un impacto social directo.

\section{Formular una hipótesis que contemple al menos dos variables y exprese de manera precisa la naturaleza y dirección de la relación entre ellas en una sola sentencia, sin explicaciones adicionales} % Numbered section

\noindent\fbox{%
	\parbox{\textwidth}{%
		\textbf{Hipótesis:} A mayor número de horas semanales trabajadas mensualmente, mayor es la apropiación y reconocimiento subjetivo de los derechos laborales por parte de las trabajadoras domésticas.
	}%
}\\

\noindent En donde las variables son:
\begin{itemize}
	\item \textbf{Independiente:} Número de horas trabajadas semanalmente de acuerdo al tipo de inserción laboral
	\item \textbf{Dependiente:} Apropiación y percepción de los derechos laborales
\end{itemize}

\section{Identificar el tipo de diseño utilizado en la investigación, especificando si es transversal o longitudinal, las fuentes de datos empleadas, los tipos de objetivos definidos y el grado de experimentación implementado, indicando si es experimental o no experimental}

\begin{itemize}
	\item \textbf{Diseño:} Es de tipo transversal, ya que se observa y analiza una situación en un úniico momento del tiempo, sin manipulación de variables.
	\item \textbf{Fuentes de datos empleadas:} Se utilizan tanto datos cuantitativos como datos cualitativos, siendo los primeros, por ejemplo, registros laborales y encuestas, y los segundos, entrevistas.
	\item \textbf{Tipos de obbjetivos definidos:}
		\begin{itemize}
			\item Descriptivos: Al explorar las caracterísitcas del grupo y las modalidades de inserción
			\item Explicativos: cuando se identifican los factores asociados a las diferencias en la apropriación de los derechos laborales
		\end{itemize}
	\item \textbf{Grado de experimentación:} No experimental, no se manipulan variables independientes ni tampoco se controlan las condiciones externas al ser una observación de fenómenos sociales tal como suceden en su entorno y contexto "natural".
\end{itemize}

%----------------------------------------------------------------------------------------

\end{document}
